\reading{MR-4}{Backtesting VaR}{grind}{4}
% ==========================================================

\begin{mrkeybox}[title={Learning objectives in this reading}]
\small
\begin{enumerate}[label=\textbf{\alph*.},leftmargin=1.7em]
\item Describe backtesting and exceptions and explain the importance of backtesting VaR models.
\item Explain the significant difficulties in backtesting a VaR model.
\item Evaluate the accuracy of a VaR model based on exceptions or failure rates by using a model
      verification test.
\item Identify and describe Type I and Type II errors in the context of a backtesting process.
\item Explain the need to consider conditional coverage in the backtesting framework.
\item Describe the Basel rules for backtesting.
\end{enumerate}
\end{mrkeybox}

% ---------------------------------------------------------
\lo{MR-4 a}{Describe backtesting and exceptions and explain the importance of backtesting VaR models.}

\begin{mrdefbox}[title={Backtesting}]
Backtesting is the verification that actual losses are reasonably consistent with projected
losses. Backtesting compares the history of VaR forecasts to actual (realized) portfolio
returns.
\end{mrdefbox}

\begin{itemize}
\item The number of observations falling outside VaR should be in line with the confidence
      level.
\item The number of exceedances is also known as the number of \term{exceptions}.
\end{itemize}

\begin{mrtrapbox}[title={The two failure directions}]
\begin{itemize}
\item \term{Too many exceptions:} the model underestimates risk.
\item \term{Too few exceptions:} excess, or inefficient, allocation of capital across units.
\end{itemize}
\end{mrtrapbox}

\begin{center}
\begin{tikzpicture}
\begin{axis}[
  width=11.2cm, height=4.4cm,
  xmin=0, xmax=60, ymin=-4.6, ymax=3.4,
  axis y line=left, axis x line=bottom, axis line style={draw=black!55},
  ylabel={\scriptsize Daily return},
  ylabel style={font=\scriptsize}, tick label style={font=\scriptsize},
  xlabel={\scriptsize Trading day}, xlabel style={font=\scriptsize},
  ytick={-4,-2,0,2}, xtick={0,10,20,30,40,50,60}, clip=false]
\addplot[FRMNavy!70, line width=0.6pt] coordinates {
(1,0.4)(2,-0.9)(3,1.1)(4,0.2)(5,-1.4)(6,0.8)(7,-0.3)(8,1.6)(9,-2.1)(10,0.5)
(11,-0.7)(12,1.2)(13,0.1)(14,-1.8)(15,2.0)(16,-0.4)(17,0.9)(18,-3.2)(19,1.4)(20,-0.6)
(21,0.3)(22,1.7)(23,-1.1)(24,0.6)(25,-2.4)(26,1.0)(27,-0.2)(28,0.8)(29,-1.6)(30,1.3)
(31,-0.5)(32,2.2)(33,-1.0)(34,0.4)(35,-3.9)(36,1.1)(37,-0.8)(38,0.7)(39,-2.0)(40,1.5)
(41,-0.1)(42,0.9)(43,-1.3)(44,2.1)(45,-0.7)(46,0.5)(47,-2.6)(48,1.2)(49,-0.9)(50,0.6)
(51,1.8)(52,-1.5)(53,0.2)(54,-3.5)(55,1.0)(56,-0.4)(57,0.8)(58,-1.9)(59,1.6)(60,-0.6)};
\addplot[FRMRed, dashed, line width=0.9pt] coordinates {(0,-3.0)(60,-3.0)};
\node[anchor=west, font=\scriptsize\sffamily, text=FRMRed] at (axis cs:60.6,-3.0)
  {$-$VaR};
\node[circle, draw=FRMRed, fill=FRMRed!25, inner sep=1.6pt] at (axis cs:18,-3.2) {};
\node[circle, draw=FRMRed, fill=FRMRed!25, inner sep=1.6pt] at (axis cs:35,-3.9) {};
\node[circle, draw=FRMRed, fill=FRMRed!25, inner sep=1.6pt] at (axis cs:54,-3.5) {};
\node[anchor=north, font=\scriptsize\sffamily, text=FRMRed] at (axis cs:35,-4.15)
  {exceptions};
\end{axis}
\end{tikzpicture}
\end{center}
\figcap{Each day's realized return is compared with the VaR forecast. The circled days are the
ones that breached it. Counting those breaches, and asking whether there are about as many as
the confidence level implies, is the whole of backtesting.}

% ---------------------------------------------------------
\lo{MR-4 b}{Explain the significant difficulties in backtesting a VaR model.}

\begin{enumerate}[label=\textbf{\alph*.}]
\item VaR models use \term{static portfolios}, but real portfolios constantly change due to
      price shifts and trading. This disparity complicates backtesting.
\item VaR models often \term{exclude various risk factors} like intraday changes and transaction
      costs, affecting the accuracy of backtesting.
\item Using \term{short time frames} for backtesting, like daily holding periods, can help
      mitigate some of the issues caused by omitted factors.
\item The backtesting sample \term{may not fully represent true risk}, making it essential to
      determine an acceptable level of exceptions.
\item Backtesting involves comparing actual returns with \term{hypothetical returns} that
      simulate VaR expectations based on static portfolios.
\item When hypothetical returns fail in backtesting, it may be necessary to \term{adjust the VaR
      modeling methodology} to enhance accuracy and risk assessment.
\end{enumerate}

% ---------------------------------------------------------
\lo{MR-4 c}{Evaluate the accuracy of a VaR model based on exceptions or failure rates by using a model verification test.}

\begin{mrdefbox}[title={Failure rate}]
An unbiased measure of the number of exceptions as a proportion of the number of samples is
called the failure rate.
\end{mrdefbox}

\begin{mrfmlbox}
\[ \text{Failure rate} \;=\; \frac{N}{T} \]
\mrvk{$N$ = number of exceptions;\quad $T$ = number of samples (observations in the backtest
window).}{}
\end{mrfmlbox}

\begin{mrexbox}[title={Example --- failure rate}]
A bank uses a 99\% confidence level VaR model over 252 trading days. The model predicts losses
exceeding the threshold on 3 days, but in reality it happens on 6 days.
\tcblower
\small
Failure rate $= 6 / 252 = \mathbf{2.38\%}$. This suggests the VaR model \term{underestimates
risk}, since exceptions occur more often than expected (1\%).
\end{mrexbox}

The number of exceptions $x$ follows a binomial probability distribution. However, we can
approximate the binomial distribution by the normal distribution.

\begin{mrfmlbox}
\[ z \;=\; \frac{x - pT}{\sqrt{p\,(1-p)\,T}} \]
\mrvk{$x$ = observed number of exceptions;\quad $p$ = the tail probability implied by the
confidence level, so $pT$ is the expected number of exceptions;\quad $T$ = number of
observations.}{}
\end{mrfmlbox}

\begin{mrexbox}[title={Example --- verifying the model}]
In 1998, daily revenue of JP Morgan fell short of the downside (95\% VaR) band on 20 days, or
more than 5\% of the time. Nine of these 20 occurrences fell within the August to October
period. Test the model over $T = 252$ days at $p = 0.05$.
\tcblower
\small
\[ z \;=\; \frac{20 - 0.05 \times 252}{\sqrt{0.05\,(1-0.05)\,252}} \;=\; \frac{7.4}{3.46}
   \;=\; 2.14 \;>\; 1.96 \]
We \term{reject} the hypothesis that the VaR model is unbiased.
\end{mrexbox}

% ---------------------------------------------------------
\lo{MR-4 d}{Identify and describe Type I and Type II errors in the context of a backtesting process.}

\begin{itemize}
\item Users of VaR models need to \term{balance Type I errors against Type II errors}.
\item Ideally, one would want to set a low Type I error rate and then have a test that creates a
      very low Type II error rate, in which case the test is said to be \term{powerful}. However,
      it is very difficult.
\end{itemize}

\begin{mrkeybox}[title={Error types}]
\begin{enumerate}[label=\textbf{\alph*.}]
\item \term{Type I error:} rejecting an accurate model.
\item \term{Type II error:} accepting an inaccurate model.
\end{enumerate}
\end{mrkeybox}

\begin{center}
\small
\begin{tabular}{>{\raggedright\arraybackslash}p{3.0cm} >{\centering\arraybackslash}p{3.4cm} >{\centering\arraybackslash}p{3.4cm}}
\arrayrulecolor{FRMRule}\toprule
\rowcolor{FRMNavy} \hdr{Decision} & \hdr{Correct} & \hdr{Incorrect}\\
Accept & OK & \textbf{Type II}\\
\rowcolor{FRMSoft} Reject & \textbf{Type I} & Power of test\\
\bottomrule
\end{tabular}
\end{center}
\figcap{Rows are the decision the backtest leads you to; columns are the true state of the
model. The two diagonal cells are the two ways of being wrong.}

\begin{center}
\begin{tikzpicture}
\begin{axis}[
  width=11.0cm, height=5.4cm,
  xmin=-3.2, xmax=8.5, ymin=0, ymax=0.47,
  xlabel={\scriptsize Number of exceptions observed},
  axis y line=none, axis x line=bottom, axis line style={draw=black!55},
  xlabel style={font=\scriptsize}, xtick={\empty},
  legend style={font=\scriptsize\sffamily, draw=FRMRule, fill=white,
    at={(0.5,-0.22)}, anchor=north, legend columns=-1, column sep=8pt},
  clip=false]
% Type I area: right tail of H0 beyond the cutoff
\addplot[domain=1.55:5.0, samples=90, smooth, draw=none, fill=FRMNavy!35, forget plot]
  {0.39894*exp(-(x^2)/2)} \closedcycle;
% Type II area: left tail of H1 below the cutoff
\addplot[domain=-0.6:1.55, samples=90, smooth, draw=none, fill=FRMRed!32, forget plot]
  {0.39894*exp(-((x-3.1)^2)/2)} \closedcycle;
\addplot[FRMNavy, line width=1.3pt, smooth, domain=-3.2:5.2, samples=170]
  {0.39894*exp(-(x^2)/2)};
\addlegendentry{Model is \textbf{accurate} ($H_0$ true)}
\addplot[FRMRed, line width=1.3pt, dashed, smooth, domain=-0.6:8.0, samples=170]
  {0.39894*exp(-((x-3.1)^2)/2)};
\addlegendentry{Model is \textbf{inaccurate} ($H_a$ true)}
\draw[black!70, line width=1.0pt] (axis cs:1.55,0) -- (axis cs:1.55,0.43);
\node[anchor=south, font=\scriptsize\sffamily, text=black!70, fill=white, inner sep=1.2pt]
  at (axis cs:1.55,0.435) {critical value};
\node[anchor=south west, font=\tiny\sffamily, text=FRMNavy, fill=white, inner sep=1.4pt,
  align=left] at (axis cs:2.35,0.26) {\textbf{Type I}\\ reject an accurate model};
\draw[FRMNavy!70, -{Stealth[length=3.5pt]}, line width=0.6pt]
  (axis cs:2.32,0.27) -- (axis cs:1.95,0.075);
\node[anchor=south east, font=\tiny\sffamily, text=FRMRed, fill=white, inner sep=1.4pt,
  align=right] at (axis cs:0.85,0.26) {\textbf{Type II}\\ accept an inaccurate model};
\draw[FRMRed!70, -{Stealth[length=3.5pt]}, line width=0.6pt]
  (axis cs:0.88,0.27) -- (axis cs:1.25,0.075);
\node[anchor=south, font=\tiny\sffamily, text=black!60, fill=white, inner sep=1.2pt]
  at (axis cs:5.6,0.16) {\textbf{power} of the test};
\end{axis}
\end{tikzpicture}
\end{center}
\figcap{The grid above says \textit{what} the two errors are; this says \textit{where} they come
from. Both errors live on the same critical value, which is why they trade off: slide the line
right and the blue Type I area shrinks while the red Type II area grows. Moving the two curves
further apart --- more observations, or a lower confidence level --- is what shrinks both at
once, and that is exactly what the improvements below propose.}

\subsection{Improvements to Type I and Type II errors}
\begin{itemize}
\item Ideally, one would want a framework that has very high \term{power}, or a high probability
      of rejecting an incorrect model.
\item The current framework could be improved by choosing a \term{lower VaR confidence level} or
      by \term{increasing the number of data observations}.
      \begin{itemize}
      \item The horizon should be \term{as short as possible} in order to increase the number of
            observations and to mitigate the effect of changes in the portfolio composition.
      \item The confidence level \term{should not be too high}, because this decreases the
            effectiveness, or power, of the statistical tests.
      \end{itemize}
\end{itemize}

\subsection{Kupiec VaR backtest}

\term{1. Unconditional coverage.} Kupiec (1995) determined a measure to accept or reject models
using the tail points of a log-likelihood ratio (LR).

The term \term{unconditional coverage} refers to the fact that we are not concerned about
independence of exception observations or the timing of when the exceptions occur.

Basically, the log-likelihood ratio $LR_{uc}$ is a test statistic to determine the hypothesis
that the model is accurate.

\begin{itemize}
\item $H_0$: accurate model
\item $H_a$: inaccurate model
\end{itemize}

\begin{mrfmlbox}
\[ LR_{uc} \;=\; -2\ln\!\left[ (1-p)^{T-N} p^{N} \right]
   \;+\; 2\ln\!\left\{ \left[ 1 - (N/T) \right]^{T-N} (N/T)^{N} \right\} \]
\mrvk{$p$ = the tail probability implied by the confidence level;\quad $N$ = number of
exceptions;\quad $T$ = number of observations;\quad $N/T$ = the observed failure rate.}{}
\end{mrfmlbox}

\begin{mrkeybox}
We would \term{reject the null hypothesis if $LR > 3.841$}.
\end{mrkeybox}

% ---------------------------------------------------------
\lo{MR-4 e}{Explain the need to consider conditional coverage in the backtesting framework.}

\begin{itemize}
\item So far the framework focuses on \term{unconditional coverage} because it ignores
      conditioning, or time variation in the data. The observed exceptions, however, could
      \term{cluster} or ``bunch'' closely in time, which also should invalidate the model.
\item In theory, these occurrences should be evenly spread over time. If, instead, we observed
      that 10 of these exceptions occurred over the last 2 weeks, this should raise a red flag.
\end{itemize}

\subsection{Conditional coverage models}

\begin{mrfmlbox}
\[ LR_{cc} \;=\; LR_{uc} \;+\; LR_{ind} \]
\mrvk{$LR_{uc}$ = the unconditional coverage statistic, which tests whether the \textit{number} of
exceptions is right;\quad $LR_{ind}$ = the independence statistic, which tests whether the
exceptions are spread out or clustered;\quad $LR_{cc}$ = the combined conditional coverage
statistic.}{}
\end{mrfmlbox}

\begin{itemize}
\item Adjust $LR_{independence}$ to remove the feature of clustering.
\item If $LR_{independence} > 3.84$, we reject independence and conclude cluster.
\item $LR_{cc} = LR_{uc} + LR_{ind}$, thus if $LR_{cc} > 5.99$, we reject the model.
\item If exceptions are determined to be serially dependent, then the model needs to be revised
      to incorporate the correlations that are evident in current conditions.
\end{itemize}

\begin{mrtrapbox}[title={Two different critical values}]
The independence test and the unconditional test are each compared against \textbf{3.84}; the
combined conditional coverage statistic is compared against \textbf{5.99}. Using the wrong one
of the two is a standard way to be caught out.
\end{mrtrapbox}

% ---------------------------------------------------------
\lo{MR-4 f}{Describe the Basel rules for backtesting.}

\begin{mrdefbox}[title={Basel Committee on backtesting}]
Basel requires that Market VaR be computed at a \term{99\% confidence interval} and backtested
for the past \term{1 year}.
\end{mrdefbox}

\begin{enumerate}[label=\textbf{\alph*.}]
\item We expect to have \term{2.5 exceptions} ($250 \times 0.01$) each year.
\item Depending on the exceptions observed, the capital multiplier (the amount a bank must hold)
      increases, and this lowers the bank's performance if the model used is incorrect.
\end{enumerate}

\begin{mrfmlbox}
\[ \text{Economic capital} \;=\; \mathrm{VaR} \times (3 + k) \]
\mrvk{$3$ = the base multiplier;\quad $k$ = the additional plus factor set by the zone the bank
falls into, so that the total multiplier runs from 3 in the green zone to 4 in the red zone.}{}
\end{mrfmlbox}

\begin{center}
\small
\begin{tabular}{>{\raggedright\arraybackslash}p{3.2cm} >{\raggedright\arraybackslash}p{4.6cm} >{\raggedright\arraybackslash}p{4.6cm}}
\arrayrulecolor{FRMRule}\toprule
\rowcolor{FRMNavy} \hdr{Zone} & \hdr{Exceptions observed} & \hdr{Multiplier}\\
\textbf{Green} & 0 to 4 observations & 3\\
\rowcolor{FRMSoft} \textbf{Yellow} & 5 to 9 observations & Increases from 3 to 4\\
\textbf{Red} & 10 or more observations & 4\\
\bottomrule
\end{tabular}
\end{center}

\subsection{The multiplier is based on the type of model error}

\begin{center}
\small
\begin{tabularx}{\linewidth}{@{}p{0.8cm} p{3.6cm} >{\raggedright\arraybackslash}X p{2.6cm}@{}}
\arrayrulecolor{FRMRule}\toprule
\rowcolor{FRMNavy} \hdr{} & \hdr{Type of error} & \hdr{Why the exceptions occurred} & \hdr{Penalty}\\
I & Basic integrity of the model is lacking & Incorrect data or errors in the programming
  & Should apply\\
\rowcolor{FRMSoft}
II & Model accuracy needs improvement & The model does not describe risks precisely
  & Should apply\\
III & Intraday trading & Positions changed during the day & Should be considered\\
\rowcolor{FRMSoft}
IV & Bad luck & Markets were particularly volatile or correlations changed; these exceptions
  should be expected to occur at least some of the time & ---\\
\bottomrule
\end{tabularx}
\end{center}

\begin{mrkeybox}
Reducing the confidence interval reduces the probability of Type II errors. Also, using longer
backtesting periods can help here.
\end{mrkeybox}


% ==========================================================
